\section{Work Done as of 24 August 2026}

\subsection{Chosen Topic}

The chosen topic is \textbf{memory consistency models} in shared-memory
multiprocessors.

\medskip
\noindent\textbf{Working title.}\\
\emph{Revisiting Sequential Consistency: Memory Consistency Models, Their Cost,
and Whether Speculation Has Made Strong Ordering Affordable}

\medskip
\noindent\textbf{Scope.}
A shared-memory multiprocessor must specify when, and in what order, one core's
writes become visible to the others.
That specification is the machine's memory consistency model, and it is an
architectural contract as binding as the instruction set.
The paper covers the design space from sequential consistency through TSO, PSO,
weak ordering and release consistency; the mechanisms that make relaxation
necessary (store buffers, out-of-order execution) and those that could make it
unnecessary (speculative execution with coherence-based detection); and the
historical argument that led the industry to reject sequential consistency
around 1990.

\medskip
\noindent\textbf{Guiding question.}
Sequential consistency was rejected as a hardware memory model on performance
grounds in the late 1980s, on machines that did not speculate.
A modern out-of-order core already contains both mechanisms a strong model
requires --- coherence invalidations that announce remote writes, and a
squash-and-restart path built for branch misprediction recovery.
The paper asks whether the performance objection still holds, and, if it does
not, what else accounts for the continued absence of sequentially consistent
hardware.

\medskip
\noindent The work is organised around four questions:

\begin{description}[leftmargin=2.2cm,style=nextline,font=\normalfont\bfseries]
\item[Q1 Formal] Can SC, TSO, PSO, weak ordering and release consistency be
  expressed as successively weaker constraints on a single set of ordering
  relations, with strict containment proved rather than asserted?
\item[Q2 Empirical] Which relaxations are \emph{observable} on real x86 and ARM
  hardware, as distinct from those merely \emph{permitted} by the architecture?
\item[Q3 Quantitative] Holding the microarchitecture fixed and varying only the
  memory model, what does sequential consistency cost with and without
  speculative load execution, as core count rises?
\item[Q4 Interpretive] If speculative SC proves inexpensive, what accounts for
  its absence from shipping processors?
\end{description}

\subsection{Mapping to the Prescribed Textbook}

The prescribed text is S.\,R.~Sarangi, \emph{Basic Computer Architecture},
version 3.09 \cite{sarangi2025basic}.
The topic maps principally to \textbf{Chapter 12, \emph{Multiprocessor
Systems}}, and within it to \textbf{\S12.4, \emph{MIMD Multiprocessors}}.

\begin{table}[H]
\centering
\small
\begin{tabular}{@{}llp{6.6cm}@{}}
\toprule
\textbf{\S} & \textbf{Title (page)} & \textbf{Role in the paper} \\
\midrule
\multicolumn{3}{@{}l}{\textbf{Primary --- Chapter 12}} \\
12.4.1 & Logical Point of View (579)          & Shared-memory abstraction. \\
12.4.2 & Coherence (581)                      & The coherence/consistency distinction, developed at length. \\
12.4.3 & \textbf{Memory Consistency (583)}    & \textbf{The core section. The paper is an extension of this.} \\
12.4.5 & Shared Caches (596)                  & Background. \\
12.4.6 & Coherent Private Caches (596)        & Extended to full MESI/MOESI and directory protocols. \\
12.4.7 & Implementing a Memory Consistency Model$^*$ (603) & Direct antecedent of the speculative-SC discussion. \\
\addlinespace
\multicolumn{3}{@{}l}{\textbf{Supporting}} \\
10.4   & Pipeline Hazards (420)               & Reordering and dependences. \\
10.9   & Performance Metrics (463)            & Benchmarking framework. \\
10.11.4& Out-of-Order Pipelines (492)         & Speculation and the squash mechanism. \\
11.2   & Caches (517)                         & Private caches, non-blocking misses. \\
11.3   & The Memory System (532)              & Latency figures motivating the store buffer. \\
12.2.2 & Shared Memory vs Message Passing (571)& Framing. \\
12.2.3 & Amdahl's Law (576)                   & Scaling discussion. \\
12.6   & Interconnection Networks (620)       & Coherence latency in the simulator. \\
App.~A & Cortex-A8/A15, Sandy Bridge (734--746)& Real-machine case studies. \\
\bottomrule
\end{tabular}
\caption{Mapping of the topic onto the prescribed text. $^*$ denotes a section
marked advanced.}
\end{table}

\noindent\textbf{Additional texts.} Because \S12.4.3 treats only sequential
consistency and the relaxed models are outside the scope of the prescribed
text, three further sources carry the technical weight of the paper.
All are fully specified, with reproduced tables of contents, in
Chapter~7 of the paper draft.

\begin{itemize}[leftmargin=*,itemsep=2pt]
\item S.\,R.~Sarangi, \emph{Next-Gen Computer Architecture: Till the End of
  Silicon}, White Falcon 2023, ISBN 8119510143 \cite{sarangi2023nextgen}.
  Chapter~9, \emph{Multicore Systems: Coherence, Consistency and Transactional
  Memory}, runs to roughly 118 pages and is the source of the formalism used
  throughout.
\item V.~Nagarajan, D.\,J.~Sorin, M.\,D.~Hill and D.\,A.~Wood, \emph{A Primer
  on Memory Consistency and Cache Coherence}, 2nd ed., Springer 2020
  \cite{nagarajan2020primer}. The standard reference on the subject.
\item Vendor architecture manuals (Intel SDM Vol.~3A; Arm ARM Ch.~B2;
  RISC-V RVWMO), cited in preference to any textbook wherever the paper states
  what a particular instruction set guarantees.
\end{itemize}

\subsection{Advisor}

I will consult \textbf{Prof.~[Sumantra Dutta Roy / Subrat Kar --- SELECT ONE]}
as advised.

\medskip
\noindent Zotero library \texttt{ELL305-MemoryConsistency} has been created and
will be shared with Prof.~Subrat Kar. Bibliography is exported to
\texttt{refs.bib} and rendered via BibTeX.

\subsection{Progress}

\subsubsection*{Summary}

\begin{table}[H]
\centering
\small
\begin{tabular}{@{}clcc@{}}
\toprule
\textbf{WP} & \textbf{Work package} & \textbf{Target} & \textbf{Complete} \\
\midrule
WP1 & Foundations: prescribed textbook          & 31 Aug & \pct{55} \\
WP2 & Advanced sources                          & 14 Sep & \pct{30} \\
WP3 & Literature survey and Zotero              & 21 Sep & \pct{35} \\
WP4 & Writing: formal, historical, methodological chapters & 24 Aug & \pct{90} \\
WP5 & Litmus experiments on real hardware       & 28 Sep & \pct{10} \\
WP6 & Simulator design and evaluation           & 26 Oct & \pct{0} \\
WP7 & Verification                              & 02 Nov & \pct{0} \\
WP8 & Final assembly and appendices             & 11 Nov & \pct{0} \\
\midrule
    & \textbf{Overall}                          &        & \pct{28} \\
\bottomrule
\end{tabular}
\caption{Work package status as of 24 August 2026.}
\end{table}

\noindent\fbox{\parbox{0.96\textwidth}{\small
\textbf{Note to the author:} every percentage in this section and in the Gantt
chart must be corrected to reflect what has \emph{actually} been done before
submission. Figures that overstate progress will be contradicted by the next
draft and are not worth the marks they might gain.}}

\subsubsection*{Gantt Chart}

\clearpage
\begin{figure}[H]
\centering
\begin{ganttchart}[
    hgrid, vgrid,
    x unit=7mm,
    y unit title=7mm,
    y unit chart=5mm,
    title height=1,
    bar height=0.6,
    group label font=\bfseries\scriptsize,
    bar label font=\scriptsize,
    milestone label font=\scriptsize\itshape,
    title label font=\scriptsize\bfseries,
    group/.append style={fill=wpcol, draw=wpcol},
    bar/.append style={fill=donecol!70, draw=donecol},
    progress label text={\scriptsize\textbf{#1\%}},
    progress label anchor/.style={left=2pt},
  ]{1}{13}

  \gantttitle{Aug}{2}
  \gantttitle{Sep}{4}
  \gantttitle{Oct}{4}
  \gantttitle{Nov}{3} \\
  \gantttitlelist{1,...,13}{1} \\

  \ganttgroup[progress=55]{WP1 Foundations}{1}{3} \\
  \ganttbar[progress=100]{T1.1 Sarangi Ch.\ 11}{1}{1} \\
  \ganttbar[progress=70]{T1.2 Sarangi \S12.1--12.4.6}{1}{2} \\
  \ganttbar[progress=20]{T1.3 Sarangi \S12.4.7}{2}{3} \\
  \ganttbar[progress=40]{T1.4 Sarangi \S10.11.4}{2}{3} \\

  \ganttgroup[progress=30]{WP2 Advanced sources}{1}{5} \\
  \ganttbar[progress=45]{T2.1 Next-Gen \S9.1--9.4}{1}{3} \\
  \ganttbar[progress=15]{T2.2 Next-Gen \S9.5--9.9}{3}{5} \\
  \ganttbar[progress=40]{T2.3 Primer Ch.\ 3--5}{1}{3} \\
  \ganttbar[progress=10]{T2.4 Primer Ch.\ 6--8, 11}{3}{5} \\

  \ganttgroup[progress=35]{WP3 Literature \& Zotero}{1}{6} \\
  \ganttbar[progress=80]{T3.1 Zotero library}{1}{2} \\
  \ganttbar[progress=30]{T3.2 Six core papers}{2}{4} \\
  \ganttbar[progress=5]{T3.3 Vendor manuals}{4}{6} \\

  \ganttgroup[progress=90]{WP4 Writing (Ch.\ 1--7)}{1}{2} \\
  \ganttbar[progress=100]{T4.1 Ch.\ 1--3}{1}{1} \\
  \ganttbar[progress=95]{T4.2 Ch.\ 4 formal}{1}{2} \\
  \ganttbar[progress=90]{T4.3 Ch.\ 5 history}{1}{2} \\
  \ganttbar[progress=85]{T4.4 Ch.\ 6 methodology}{1}{2} \\

  \ganttgroup[progress=10]{WP5 Litmus experiments}{2}{7} \\
  \ganttbar[progress=25]{T5.1 Test harness}{2}{3} \\
  \ganttbar[progress=5]{T5.2 Test catalogue}{3}{5} \\
  \ganttbar[progress=0]{T5.3 Run on x86, AArch64}{5}{6} \\
  \ganttbar[progress=0]{T5.4 Write Ch.\ 8}{6}{7} \\

  \ganttgroup[progress=0]{WP6 Simulator}{4}{11} \\
  \ganttbar[progress=0]{T6.1 Cores, store buffers}{4}{6} \\
  \ganttbar[progress=0]{T6.2 MESI directory}{5}{7} \\
  \ganttbar[progress=0]{T6.3 Ordering policies}{7}{9} \\
  \ganttbar[progress=0]{T6.4 Workloads, sweeps}{8}{10} \\
  \ganttbar[progress=0]{T6.5 Write Ch.\ 9--10}{9}{11} \\

  \ganttgroup[progress=0]{WP7 Verification}{9}{12} \\
  \ganttbar[progress=0]{T7.1 Axiomatic checker}{9}{11} \\
  \ganttbar[progress=0]{T7.2 Write Ch.\ 11}{11}{12} \\

  \ganttgroup[progress=0]{WP8 Final assembly}{10}{13} \\
  \ganttbar[progress=0]{T8.1 Ch.\ 12--13}{10}{12} \\
  \ganttbar[progress=0]{T8.2 Appendices A--G}{11}{13} \\
  \ganttbar[progress=0]{T8.3 Revision passes}{12}{13} \\

  \ganttmilestone{v0.1 (24 Aug)}{2} \\
  \ganttmilestone{v0.2}{7} \\
  \ganttmilestone{v0.3}{11} \\
  \ganttmilestone{v1.0 (11 Nov)}{13}
\end{ganttchart}
\caption{Project schedule. Week 1 commences 17 August 2026; week 13 commences
9 November 2026. Shaded portions of each bar indicate completion. Milestone
dates for v0.2 and v0.3 are estimates pending confirmation.}
\end{figure}
\clearpage

\subsubsection*{Work Package Detail}

\paragraph{WP1 --- Foundations (\pct{55}).}
\begin{itemize}[leftmargin=*,itemsep=1pt]
\item \textbf{T1.1} Sarangi \cite{sarangi2025basic} Ch.~11, \emph{The Memory
  System}, \S11.1--11.3 (pp.~505--560). \pct{100}. Cache organisation, locality,
  and the latency figures that motivate the store buffer.
\item \textbf{T1.2} Sarangi Ch.~12, \S12.1--12.4.6 (pp.~565--603). \pct{70}.
  \S12.4.2 and \S12.4.3 read closely; \S12.4.6 in progress.
\item \textbf{T1.3} Sarangi \S12.4.7, \emph{Implementing a Memory Consistency
  Model} (pp.~603--608). \pct{20}. Begun; requires \S12.4.6 first.
\item \textbf{T1.4} Sarangi \S10.11.4, \emph{Out-of-Order Pipelines}
  (pp.~492--505). \pct{40}. Needed for the squash mechanism central to Q3.
\end{itemize}

\paragraph{WP2 --- Advanced sources (\pct{30}).}
\begin{itemize}[leftmargin=*,itemsep=1pt]
\item \textbf{T2.1} Sarangi \cite{sarangi2023nextgen} \S9.1--9.4. \pct{45}.
  \S9.4 supplies the execution-witness and access-graph notation adopted
  throughout the paper, chosen over the alternatives for continuity with the
  course.
\item \textbf{T2.2} Sarangi \S9.5--9.9 (coherence protocols, memory models,
  race detection, transactional memory). \pct{15}.
\item \textbf{T2.3} Nagarajan et al.\ \cite{nagarajan2020primer} Ch.~3--5
  (SC, TSO, relaxed models). \pct{40}.
\item \textbf{T2.4} Nagarajan et al.\ Ch.~6--8, 11 (coherence protocols,
  specification and validation). \pct{10}.
\end{itemize}

\paragraph{WP3 --- Literature and Zotero (\pct{35}).}
\begin{itemize}[leftmargin=*,itemsep=1pt]
\item \textbf{T3.1} Zotero library created; 22 entries filed with PDFs.
  \pct{80}. Exported to \texttt{refs.bib}, under version control.
\item \textbf{T3.2} Six papers identified as load-bearing for the argument and
  to be read in full rather than consulted: Lamport
  \cite{lamport1979multiprocessor}; Dubois, Scheurich and Briggs
  \cite{dubois1986memory}; Adve and Hill \cite{adve1990weak}; Adve and
  Gharachorloo \cite{adve1996shared}; Gniady, Falsafi and Vijaykumar
  \cite{gniady1999sc}; Alglave, Maranget and Tautschnig
  \cite{alglave2014herding}. \pct{30}.
\item \textbf{T3.3} Vendor manuals \cite{intel_sdm,arm_arm,riscv_spec}.
  \pct{5}.
\end{itemize}

\noindent Further references filed and cited: Gharachorloo et al.\
\cite{gharachorloo1990memory}; Hill \cite{hill1998multiprocessors}; Manson,
Pugh and Adve \cite{manson2005java}; Boehm and Adve
\cite{boehm2008foundations}; Ceze et al.\ \cite{ceze2007bulksc}; Blundell,
Martin and Wenisch \cite{blundell2009invisifence}; Marino et al.\
\cite{marino2011sc}. Supporting texts: Hennessy and Patterson
\cite{hennessy2019quantitative}; Culler, Singh and Gupta
\cite{culler1998parallel}; Herlihy et al.\ \cite{herlihy2020art}.

\paragraph{WP4 --- Writing (\pct{90}).}
A 42-page draft of the paper proper exists in the repository
(\texttt{main.tex}), comprising:

\begin{table}[H]
\centering
\small
\begin{tabular}{@{}rlc@{}}
\toprule
\textbf{Ch.} & \textbf{Title} & \textbf{State} \\
\midrule
1 & Introduction --- motivating example, problem statement, contributions & drafted \\
2 & Background --- memory wall, caches, store buffer, out-of-order & drafted \\
3 & Cache Coherence, and What It Does Not Solve & drafted \\
4 & The Design Space of Memory Consistency Models & drafted \\
5 & Historical Evolution & drafted \\
6 & Methodology & drafted \\
7 & Source Texts and Their Coverage & drafted \\
\bottomrule
\end{tabular}
\caption{State of the paper draft. Seven chapters, 42 pages, five TikZ figures,
compiling without errors or undefined references.}
\end{table}

\noindent All figures are drawn in TikZ per the course requirement. The
document is built with \texttt{latexmk}; \texttt{main.pdf} is deliberately
excluded from version control and attached instead to tagged releases.

\paragraph{WP5 --- Litmus experiments (\pct{10}).}
Design complete and documented in Chapter~6 of the paper draft; implementation
begun. The harness must avoid \texttt{pthread\_barrier}, whose internal fences
suppress the window under observation, in favour of a sense-reversing spin
barrier built from relaxed atomics. Shared variables are padded to distinct
cache lines to prevent false sharing from confounding the coherence behaviour.

\paragraph{WP6--WP8 (\pct{0}).}
Not started; scheduled per the Gantt chart. Designs for WP6 (simulator) and
WP7 (verification) are specified in Chapter~6 of the paper draft, stated in
advance so that the methodology can be assessed independently of the results.

\subsection{Risks and Mitigations}

\begin{table}[H]
\centering
\small
\begin{tabular}{@{}p{4.6cm}cp{7.0cm}@{}}
\toprule
\textbf{Risk} & \textbf{Sev.} & \textbf{Mitigation} \\
\midrule
Litmus tests yield zero observations, making Ch.~8 empty
 & Med & Vary thread placement, optimisation level and inserted delay; a null
        result with the parameter sweep documented is itself reportable. \\
Simulator scope creep consumes the schedule
 & High & Feature list frozen after v0.2. Purpose-built simulator is the
         baseline; Tejas \cite{sarangi2023nextgen} evaluated in parallel as a
         stretch goal only, never a dependency. \\
No AArch64 hardware available for Q2
 & Med & Single-board computer, Android device under Termux, or institutional
         server; failing all three, Q2 narrows to x86 and the limitation is
         stated. \\
Topic requires material beyond the course
 & Low & Sarangi's companion volume \cite{sarangi2023nextgen} is pitched at
        exactly this level and is by the same author. \\
\bottomrule
\end{tabular}
\end{table}

\subsection{Plan to v0.2}

\begin{enumerate}[leftmargin=*,itemsep=2pt]
\item Complete WP1 (T1.2--T1.4) and advance WP2 to \pct{70}.
\item Read the six load-bearing papers in full (T3.2 to \pct{100}).
\item Implement the litmus harness and the full test catalogue; obtain
  observation counts on x86 and, if hardware permits, AArch64.
\item Draft Chapter~8 (Results I) from those measurements.
\item Begin WP6: cores, store buffers and load queue.
\item Insert the measured SB violation count into \S1.1 of the paper, so that
  the opening example is a measurement rather than an assertion.
\end{enumerate}

\subsection{Repository}

All work is under version control at\\
\texttt{https://github.com/y4shv/ELL305-term-paper} (private; access granted to
the instructor on request).

\medskip
\noindent Draft \#1 corresponds to tag \texttt{v0.1}.
\texttt{draft1.tex} builds this progress report; \texttt{main.tex} builds the
paper draft described in WP4.
\texttt{refs.bib} is exported from Zotero.
